% =============================================================================
% OpenScent — draft v1, 2026-09-09
% Set THIS file as the main document in Overleaf (not template.tex).
% Content source: openscent/reports/paper-draft-v1.md
% Every figure re-derived from pipeline/status.py on 2026-09-09 — re-run before
% submission, several moved within the last week.
% =============================================================================
% `chemistry` is deliberately NOT enabled. It loads mhchem, chemfig, chemnum and
% chemmacros (cls lines 514-517), and this paper uses none of them — no \ce{}, no
% \compound{}, no schemes. Enabling it roughly tripled compile time and put the
% project near Overleaf's free-plan timeout. Add it back the day a structure
% diagram or a reaction actually appears in the text, and not before.
\documentclass[
  computerscience
]{genesisl1-paper}

% -----------------------------------------------------------------------------
% Document identity
% -----------------------------------------------------------------------------
\GenesisTitle{OpenScent: a permissively licensed structure--odour corpus}
\GenesisSubtitle{Built by construction rather than by permission}
\GenesisRunningTitle{OpenScent: a CC0 structure--odour corpus}
\GenesisVersion{Draft 0.1}
\GenesisStatus{DRAFT}
\GenesisDate{\today}
\GenesisAuthors{Ivan Pezzini}                       % TODO authorship — see notes
\GenesisAffiliations{Independent researcher. Project proposed by Mikhail Fedorov, GenesisL1.}
\GenesisCorrespondence{veryveryarte@gmail.com}
\GenesisIdentifier{}
\GenesisRepository{https://github.com/VanPez/openscent}
\GenesisFields{Cheminformatics; scientific data; olfaction}
\GenesisKeywords{structure--odour relationship; open data; CC0; text mining; provenance}
\GenesisLicense{Data CC0 1.0; code Apache-2.0}

% Optional exact-artifact metadata. Fill at submission, not before.
\GenesisDataIdentifier{}
\GenesisCodeIdentifier{}
\GenesisModelIdentifier{}
\GenesisExecutionIdentifier{}
\GenesisContentHash{}
\GenesisIPFSCID{}
\GenesisIPFSGateway{}

\begin{document}

\begin{GenesisFrontMatter}
  \begin{GenesisStructuredAbstract}
    \AbstractBackground{
      Machine learning on olfaction depends on datasets pairing molecular structures
      with odour descriptors. The descriptor-labelled sets in common use derive from
      commercial flavour-and-fragrance databases carrying non-commercial licence terms,
      which foreclose commercial application, tokenised derivatives and independent
      redistribution. The obstacle is licensing rather than science.
    }
    \AbstractMethods{
      Odour assertions were extracted from two public-domain sources: 5{,}346 United
      States patent documents, which are government works not subject to copyright, and
      1{,}018 Hazardous Substances Data Bank records obtained through PubChem. Every
      molecule name and descriptor was required to appear as a literal substring of its
      source sentence, an invariant re-asserted on every run. A human adjudicated every
      row against written rules; automated proposals were advisory and never
      auto-accepted. Vocabulary admission required attestation across at least twenty
      distinct documents, counted at passage level.
    }
    \AbstractResults{
      The corpus contained 1{,}193 distinct molecules across a 67-term vocabulary, of
      which 20 terms reached the threshold of 30 molecules with adjudication ongoing
      (418 productive rows outstanding, a lower bound, targeting the terms nearest threshold). A
      648-molecule subset
      carrying 22 computed structural descriptors supported supervised prediction of
      odour labels under five-fold stratified cross-validation, with areas under the
      ROC curve of 0.93 for \emph{fatty} ($n=16$) and 0.87 for \emph{mint} ($n=24$).
      Four negative results constrained the method: per-compound note tiers were absent
      from 60 patents examined; 29 additional vocabulary candidates were refused for
      insufficient molecular support; the pre-1930 literature was redundant with
      existing sources for 27 of 30 compounds tested; and extraction at sentence scope
      was identified as the binding recall constraint.
    }
    \AbstractConclusion{
      A permissively licensed structure--odour corpus can be assembled by construction
      from public-domain sources. The result is smaller than the encumbered
      alternatives and is reported as such, and its labels carry structure-dependent
      signal rather than noise.
    }
  \end{GenesisStructuredAbstract}
\end{GenesisFrontMatter}

% =============================================================================
\section{Introduction}

Quantitative structure--odour relationship work requires training data mapping molecular
structures to odour descriptors. The dataset most widely used for this purpose comprises
3{,}523 molecules with expert-assigned descriptors, curated from the commercial
\emph{Leffingwell PMP 2001} database and released alongside the reference deep-learning
treatment of the problem \cite{sanchez2019machine}. It is distributed under CC~BY-NC, and
its files are additionally access-restricted, released only on request and on condition
of non-commercial use \cite{leffingwell2020}.

The consequence is not scientific but legal. A non-commercial clause forecloses
commercial application, forecloses tokenised or tradeable derivatives, and prevents a
third party from redistributing a corrected or extended version. Work built upon such
data cannot in turn be built upon.

Permissively licensed alternatives exist but differ in kind and in scale. Keller and
Vosshall \cite{keller2016} released psychophysical ratings for 480 molecules under CC0:
panel ratings against a fixed set of semantic scales rather than curated descriptor
assignments. The resource is valuable and does not substitute for a descriptor-labelled
corpus.

We report an attempt to obtain the missing artefact by construction, rebuilding a
molecule-to-descriptor corpus from sources carrying no restriction, so that the licence
question is settled before any row exists rather than negotiated afterwards.

Two source classes qualify. United States patent full text is a government work not
subject to copyright, and is where modern aroma chemicals are described, including
captive molecules documented nowhere else. The Hazardous Substances Data Bank, accessed
through PubChem, is public domain and supplies a compound identifier, making the link
from an odour claim to a structure exact rather than inferred from a name.

\subsection{Aims}
\begin{enumerate}
  \item Determine whether a descriptor-labelled odour corpus of usable size can be
        assembled entirely from public-domain sources.
  \item Establish extraction and adjudication procedures under which every recorded
        claim remains traceable to a verbatim source span.
  \item Evaluate whether the resulting labels carry structure-dependent signal.
\end{enumerate}

% =============================================================================
\section{Methods}

Four procedural rules governed the work. Each is stated as a rule and as the failure it
prevents, because in every case the failure occurred first.

\subsection{Subjects, materials, and datasets}

Patent text was harvested from 5{,}346 United States documents identified through
classification searches of C11B~9/00 (essential oils and perfumes) and A61Q~13/00
(perfume formulations). Non-US patents were excluded throughout: they do not carry the
absence of copyright on which the licence claim depends. The European Patent Office Open
Patent Services interface was used to \emph{identify} United States documents; only
United States documents were retrieved.

Odour annotations were additionally obtained from the Hazardous Substances Data Bank
through PubChem, yielding 1{,}018 usable records across 651 distinct compound
identifiers after six records were retired on provenance or attribution grounds. The
database is a hazardous-substances resource, so industrial compounds appear alongside
aroma chemicals.

Sources derived from GoodScents or Leffingwell were excluded from every stage.

\subsection{Extraction procedure}

\textbf{Verbatim extraction.} Every molecule name and descriptor recorded in a row was
required to be a literal substring of its source sentence. A status utility re-asserted
this invariant across the corpus on each run and refused to report figures when it was
violated. A paraphrased or generated row is indistinguishable from an extracted one once
written; per-row provenance is meaningful only when the recorded span can be located in
the cited source. Text normalisation was consequently versioned and applied identically
to source text and extracted spans, a comparison between differently normalised strings
not constituting a verbatim check.

\textbf{Yield measurement before retrieval.} Before harvesting a classification, a
200-document sample was retrieved and extraction yield measured, at a cost of
approximately twenty minutes. This altered the decision three times
(Table~\ref{tab:yield}).

\begin{table}[t]
\caption{Sampled yield and the resulting retrieval decision.}
\label{tab:yield}
\centering
\begin{tabular}{lcl}
\hline
Classification & Candidates/document & Decision \\
\hline
C11B~9/00 (subtree) & 1.02 & retrieved \\
A23L~27/00          & 0.17 & declined \\
C11D~3/50           & 0.36 & declined \\
\hline
\end{tabular}
\end{table}

Yield alone proved insufficient. C11D~3/50 (detergent perfumery) was sampled because it
was expected to carry amber and musk descriptors, and yielded twice the rate of the
declined food classification. Counting where its descriptors fell showed approximately
70\,\% landing on vocabulary terms already well past threshold, while \emph{sandalwood}
and \emph{animalic}, the two terms motivating the sample, did not occur. Target coverage
is now measured alongside yield.

\textbf{Adjudicated review.} A human decided every row against written rules; automated
proposals were advisory and automatic rejection was disabled. Blind proposal accuracy,
measured by withholding proposals on 30 rows, was 67\,\%; errors skewed towards rejecting
valid rows. The rules were written after a blind evaluation disagreed on 14 of 98 rows
and the disagreements proved to be two unwritten conventions applied in opposite
directions. Three rules required the most deliberation: mixtures yield no row, including
mixtures of stereoisomers of a single compound, since enantiomers may smell differently
and the odour of a racemate is not that of either component; a compound named as a
comparison is a reference point rather than a descriptor, whereas a plant or food named
as a smell is a descriptor; and negation or comparison inverts the claim.

\textbf{Independent attestation.} A vocabulary term was admitted only on appearing in at
least 20 distinct documents with at least 30 attestations, an attestation being a
distinct sentence. Four successive counters were required, each correcting the
predecessor's blind spot: raw occurrences, document counts, a per-term boilerplate ratio,
and finally passage-level counting. The admitted vocabulary comprises 67 terms with 119
surface forms; the mapping from surface form to term is a human-written table and is the
only point at which extracted text becomes a category.

\subsection{Statistical analysis}

Structural descriptors were computed offline with RDKit from structures resolved through
PubChem identifiers: 22 interpretable descriptors comprising molecular weight, log~$P$,
topological polar surface area, hydrogen-bond donors and acceptors, rotatable bonds, ring
counts, fraction sp\textsuperscript{3}, element counts, Labute surface area and Bertz
complexity. Descriptors were computed rather than retrieved; requesting precomputed
properties from a remote service would make the feature table depend on that service's
descriptor version, which is neither reproducible nor auditable.

One binary classification task was defined per label. Gradient-boosted decision trees were
evaluated under five-fold stratified cross-validation, reporting area under the ROC curve.
Labels with fewer than ten positive instances were excluded.

\begin{GenesisProvenance}
Corpus, ontology, extraction code and analysis scripts are at
\url{https://github.com/VanPez/openscent}. Every row records its source document
identifier, the verbatim source sentence, the extracted spans and a licence basis.
All headline figures are computed by a single utility so that any quoted number can be
regenerated rather than recalled.
\end{GenesisProvenance}

% =============================================================================
\section{Results}

\subsection{Corpus}

Harvesting produced 5{,}346 patent documents, of which 2{,}426 were represented in the
review queue as 6{,}686 candidate rows. Of these, 955 were adjudicated by hand (627
accepted, 326 rejected, 2 deferred). Together with 1{,}018 database records the corpus
contained 1{,}193 distinct molecules.

Twenty of the 67 vocabulary terms reached the threshold of 30 molecules:
\emph{acid, aromatic, citrus, earthy, fatty, floral, fresh, fruity, green, herbal, lily,
mint, musk, musty, powdery, pungent, rose, spicy, sweet} and \emph{woody}. Five further
terms lay within seven molecules of threshold: \emph{sandalwood} (28), \emph{amber} (28),
\emph{camphoraceous} (27), \emph{animalic} (23) and \emph{aldehydic} (23).

The design target was 60--100 terms at 30 molecules each. The corpus reached 20 at the
census date, with adjudication incomplete: 418 of the 5{,}731 undecided rows satisfy the
review queue's productivity criterion, carrying both a compound-like token and a
vocabulary term, and they fall disproportionately on terms below threshold
(\emph{aldehydic} 30 rows, \emph{camphoraceous} 16, \emph{sandalwood} 14, \emph{animalic}
9). That count is a lower bound: the review interface tests against a hard-coded
vocabulary of 110 surface forms while the ontology now holds 119, so rows carrying only
one of the 46 unlisted forms are not counted. At the precision observed within the
productive subset (0.82), these adjudications alone would be expected to carry several
further terms past threshold without any additional source material. Twenty is therefore
a census figure, not a ceiling.

\subsection{Structural signal}

Of 651 database compounds, 648 resolved to parseable structures; three failures were
hypervalent halogen oxidisers rejected by the valence model. Cross-validated performance
is given in Table~\ref{tab:auc}.

\begin{table}[t]
\caption{Five-fold stratified cross-validated AUC by label, with positive counts.}
\label{tab:auc}
\centering
\begin{tabular}{lrc}
\hline
Label & $n$ & AUC \\
\hline
chlorine  &  18 & 0.943 \\
fatty     &  16 & 0.925 \\
mint      &  24 & 0.874 \\
rose      &  12 & 0.873 \\
citrus    &  11 & 0.862 \\
aromatic  & 140 & 0.851 \\
fruity    &  85 & 0.843 \\
floral    &  39 & 0.820 \\
green     &  12 & 0.799 \\
pungent   & 186 & 0.751 \\
woody     &  10 & 0.739 \\
musty     &  28 & 0.692 \\
spicy     &  11 & 0.657 \\
acid      &  31 & 0.633 \\
sweet     & 149 & 0.619 \\
\hline
\end{tabular}
\end{table}

\begin{GenesisEvidence}
Three qualifications belong with Table~\ref{tab:auc} and should not be separated from it.
The result for \emph{chlorine} is determined by halogen count and serves as a pipeline
check rather than a finding. Labels with 10--12 positive instances yield unstable
cross-validated estimates at this sample size. The label \emph{aromatic} carries a
semantic confound: the odour term and the structural property are distinct concepts that
correlate, so part of that score is unearned. The substantive results are \emph{fatty}
(0.925) and \emph{mint} (0.874), where chain length and lipophilicity track fattiness and
mint labels concentrate within a narrow structural family.
\end{GenesisEvidence}

The purpose of this analysis was negative control rather than performance. A corpus of
incorrect labels, however carefully assembled, would score near 0.5 and would be
indistinguishable from a correct one by inspection.

% =============================================================================
\section{Discussion}

The study shows that a descriptor-labelled odour corpus can be assembled entirely from
public-domain sources, and that the resulting labels carry structure-dependent signal.
The corpus is smaller than encumbered alternatives; its distinguishing property is that
it may be used commercially, extended, and redistributed by anyone.

The transferable content of this work is procedural, together with four results that
constrain what the sources can be asked to supply.

\textbf{Note tiers are absent from patent text.} Perfumery classifies materials as top,
heart or base notes. Across 60 patents examined there were no per-compound tier
assertions: the terms describe accord architecture, never an individual molecule. A rule
fitted to curated tier assignments reproduced them at 56\,\% under leave-one-out
validation, marginally above chance across three classes. Tiers are neither sourceable
from this literature nor recoverable by inference.

\textbf{Vocabulary was not the binding constraint.} Twenty-nine further descriptor
candidates passed an independence test and were nonetheless refused: the best supported
reached 8 molecules against a threshold of 30. Admitting them would have converted 20 of
67 into 20 of 95, leaving the numerator unchanged and the denominator worse.

\textbf{The pre-1930 literature was redundant rather than inconvenient.} Of 30 classic
synthetics from that period, 27 already carried a database odour record; apparent misses
were naming artefacts. Expired-copyright sources were dropped for redundancy, not for
licensing.

\textbf{Sentence scope is the binding recall constraint.} Three independent observations
converge: a heading-matching rule was removed for producing rows naming no molecule; 345
sentences carry descriptors whose compound is named in the preceding sentence; and the
single standing recall failure in the evaluation set is of that form. Extraction at
passage scope is the identified next step and is not yet implemented.

\begin{GenesisInterpretation}
Our results provide evidence that the licensing obstacle in structure--odour data is
tractable by source selection rather than by negotiation. The observation that patent
text supplies descriptors but not note tiers points towards a general limit: patent
prose documents what a compound smells of, not how a perfumer would deploy it. These
findings are consistent with the corpus remaining useful for descriptor prediction while
being unsuitable, without further sources, for tasks requiring formulation semantics.
\end{GenesisInterpretation}

\begin{GenesisLimitations}
Although encouraging, this study has limitations. Twenty of 67 vocabulary terms reached
threshold against a design target of 60--100, with at least 418 productive adjudications
outstanding. The patent-derived half of the corpus does
not yet carry computed structures, so the analysis in Table~\ref{tab:auc} covers 648 of
1{,}193 molecules. Molecules are matched between the two sources by name text rather than
by structure, so stereochemical variants remain distinct entries. Whole-queue precision
of the candidate filter is approximately 0.2; a figure of 0.82 recorded elsewhere in the
project is precision within the productive subset and must not be read as the filter's
precision. A text-containment audit identified 835 groups covering 1{,}035 documents,
19.4\,\%, as the same disclosure published more than once, typically an application and
its granted patent; this is an upper bound, since the detector admits era-specific
template text as evidence. Redundancy does not inflate the molecule counts, which are
computed over sets, but it bears on the document-count criterion for vocabulary
admission, and that recount has not been performed. Despite these limitations, the
consistency of the observed structural signal across independently adjudicated labels
supports the corpus being fit for its stated purpose.
\end{GenesisLimitations}

% =============================================================================
\section{Conclusion}

A structure--odour corpus of 1{,}193 molecules across 20 vocabulary terms was assembled
entirely from public-domain sources and released without restriction. Its labels carry
structure-dependent signal. The corpus is smaller than licence-encumbered alternatives,
and unlike them it may be used, extended and redistributed by anyone.

% =============================================================================
\section*{Data and code availability}

Corpus and ontology are released under CC0~1.0; code under Apache-2.0. Repository:
\url{https://github.com/VanPez/openscent}. Every row carries its source identifier, the
verbatim source sentence and a licence basis, so that a challenge to any single source
removes the affected rows rather than invalidating the corpus.

CC0 waives copyright, not patents. The corpus is built from patent text; the presence of
a molecule in it says nothing about live patent claims covering that molecule or its use.

% \PrintGenesisArtifactRecord

\section*{Author contributions}
% TODO CRediT statement. Mikhail Fedorov proposed the project; authorship unsettled.

\section*{Ethics, privacy, rights, and safety}

The work uses only published patent text and public-domain database records. No personal
data, credentials or confidential material are present in the corpus or in any published
artefact.

\section*{Competing interests and funding}
% TODO declare GenesisL1 relationship.

\section*{Acknowledgements}
% TODO

\GenesisEndMatterNotice

\bibliographystyle{plainnat}
\bibliography{references}

\end{document}
